\textbf{Phần 1. Trắc nghiệm lựa chọn}
\setcounterpageref{ex}{0}

\begin{ex} %[5D1N1]
Số gồm $3$ chục triệu, $8$ nghìn, $2$ chục và $1$ đơn vị được viết là
\choice
{$30\,800\,021$}
{$3\,008\,021$}
{$30\,080\,021$}
{\True $30\,008\,021$}
\loigiai{
Số gồm $3$ chục triệu ($30\,000\,000$), $8$ nghìn ($8\,000$), $2$ chục ($20$) và $1$ đơn vị ($1$) được viết là $30\,008\,021$.\\
\textbf{Đáp án của câu hỏi là D}
}
\end{ex}

\begin{ex} %[5D1N1]
Số $52\,814$ được đọc là
\choice
{\True Năm mươi hai nghìn tám trăm mười bốn}
{Năm mươi hai nghìn tám trăm mười}
{Năm trăm hai mươi tám nghìn mười bốn}
{Năm mươi hai triệu tám trăm mười bốn}
\loigiai{
Số $52\,814$ có $5$ chục nghìn, $2$ nghìn, $8$ trăm, $1$ chục và $4$ đơn vị, được đọc là: Năm mươi hai nghìn tám trăm mười bốn.\\
\textbf{Đáp án của câu hỏi là A}
}
\end{ex}

\begin{ex} %[5D1H1]
Số tự nhiên gồm $8$ trăm nghìn, $2$ chục nghìn, $1$ chục và $5$ đơn vị là
\choice
{$82\,015$}
{\True $820\,015$}
{$802\,015$}
{$820\,150$}
\loigiai{
Số gồm $8$ trăm nghìn, $2$ chục nghìn, $0$ nghìn, $0$ trăm, $1$ chục và $5$ đơn vị là $820\,015$.\\
\textbf{Đáp án của câu hỏi là B}
}
\end{ex}

\begin{ex} %[5D1H1]
Chữ số $7$ trong số $863\,749$ thuộc hàng nào và có giá trị bằng bao nhiêu?
\choice
{Hàng chục, giá trị là $70$}
{Hàng nghìn, giá trị là $7\,000$}
{\True Hàng trăm, giá trị là $700$}
{Hàng chục nghìn, giá trị là $70\,000$}
\loigiai{
Trong số $863\,749$, kể từ phải sang trái: chữ số $9$ thuộc hàng đơn vị, chữ số $4$ thuộc hàng chục, chữ số $7$ thuộc hàng trăm nên có giá trị là $700$.\\
\textbf{Đáp án của câu hỏi là C}
}
\end{ex}

\begin{ex} %[5D1H1]
Số thích hợp điền vào dấu hỏi chấm để $504\,842 = 500\,000 + 4\,000 + ? + 40 + 2$ là
\choice
{$80$}
{$8\,000$}
{$8$}
{\True $800$}
\loigiai{
Phân tích số $504\,842$ theo các hàng:
\[
504\,842 = 500\,000 + 4\,000 + 800 + 40 + 2.
\]
Vậy số cần điền là $800$.\\
\textbf{Đáp án của câu hỏi là D}
}
\end{ex}

\begin{ex} %[5D1H1]
Số tự nhiên $1\,200\,324$ được phân tích thành tổng các hàng là
\choice
{$1\,000\,000 + 20\,000 + 300 + 20 + 4$}
{\True $1\,000\,000 + 200\,000 + 300 + 20 + 4$}
{$1\,000\,000 + 200\,000 + 3\,000 + 20 + 4$}
{$1\,200\,000 + 300 + 24$}
\loigiai{
Số $1\,200\,324$ gồm $1$ triệu, $2$ trăm nghìn, $3$ trăm, $2$ chục và $4$ đơn vị:
\[
1\,200\,324 = 1\,000\,000 + 200\,000 + 300 + 20 + 4.
\]
\textbf{Đáp án của câu hỏi là B}
}
\end{ex}

\begin{ex} %[5D1H1]
Ba số chẵn liên tiếp được viết vào $3$ chiếc mũ của Việt, Nam và Rô-bốt. Rô-bốt nhìn thấy số viết trên mũ của Việt và Nam lần lượt là $2\,032$ và $2\,028$. Số trên mũ của Rô-bốt là
\choice
{\True $2\,030$}
{$2\,029$}
{$2\,034$}
{$2\,026$}
\loigiai{
Ba số chẵn liên tiếp cách nhau $2$ đơn vị. Số nằm giữa $2\,028$ và $2\,032$ là $2\,030$. Vậy chiếc mũ Rô-bốt đang đội mang số $2\,030$.\\
\textbf{Đáp án của câu hỏi là A}
}
\end{ex}

\begin{ex} %[5D1H1]
Cho độ cao của ba đỉnh núi: Đỉnh E-vơ-rét ($8\,848\text{ m}$), Đỉnh núi Bà Đen ($986\text{ m}$), Đỉnh Ô-lim-pót Mon ($25\,000\text{ m}$). Đỉnh núi nào cao nhất?
\choice
{Đỉnh núi Bà Đen}
{Đỉnh E-vơ-rét}
{Cả ba đỉnh cao bằng nhau}
{\True Đỉnh Ô-lim-pót Mon}
\loigiai{
So sánh các số: $986 < 8\,848 < 25\,000$. Vậy đỉnh núi cao nhất là đỉnh Ô-lim-pót Mon ($25\,000\text{ m}$).\\
\textbf{Đáp án của câu hỏi là D}
}
\end{ex}

\begin{ex} %[5D1H1]
Trong ba bộ đồ chơi: Ô tô ($195\,000$ đồng), Ngôi nhà búp bê ($1\,175\,000$ đồng), Máy chơi game ($2\,000\,000$ đồng). Bộ đồ chơi nào có giá tiền thấp nhất?
\choice
{\True Ô tô}
{Ngôi nhà búp bê}
{Máy chơi game}
{Ô tô và Ngôi nhà búp bê}
\loigiai{
So sánh giá tiền: $195\,000 < 1\,175\,000 < 2\,000\,000$. Vậy bộ đồ chơi có giá thấp nhất là ô tô ($195\,000$ đồng).\\
\textbf{Đáp án của câu hỏi là A}
}
\end{ex}

\begin{ex} %[5D1H1]
Số sản phẩm bán được của một cửa hàng trong bốn năm là: Năm 2020 bán $2\,873$ sản phẩm; Năm 2021 bán $2\,837$ sản phẩm; Năm 2022 bán $3\,293$ sản phẩm; Năm 2023 bán $3\,018$ sản phẩm. Năm cửa hàng bán được nhiều sản phẩm nhất là
\choice
{Năm 2020}
{Năm 2021}
{\True Năm 2022}
{Năm 2023}
\loigiai{
So sánh số sản phẩm: $2\,837 < 2\,873 < 3\,018 < 3\,293$. Vậy năm $2022$ bán được nhiều sản phẩm nhất ($3\,293$ sản phẩm).\\
\textbf{Đáp án của câu hỏi là C}
}
\end{ex}

\begin{ex} %[5D1H1]
Doanh thu thực tế trong một ngày của một cửa hàng là $2\,545\,000$ đồng. Cuối ngày, người bán hàng nói: ``Doanh thu cả ngày của cửa hàng được khoảng $2\,500\,000$ đồng''. Người bán hàng đã làm tròn số tiền đến hàng nào?
\choice
{Hàng chục nghìn}
{\True Hàng trăm nghìn}
{Hàng triệu}
{Hàng nghìn}
\loigiai{
Số $2\,545\,000$ có chữ số hàng chục nghìn là $4 < 5$, khi làm tròn đến hàng trăm nghìn ta giữ nguyên chữ số hàng trăm nghìn là $5$ và các chữ số phía sau chuyển thành $0$, được $2\,500\,000$.\\
\textbf{Đáp án của câu hỏi là B}
}
\end{ex}

\begin{ex} %[5D1V1]
Khi làm tròn số $2\,545\,000$ đến hàng chục nghìn, ta được số nào sau đây?
\choice
{$2\,540\,000$}
{$2\,500\,000$}
{\True $2\,550\,000$}
{$2\,600\,000$}
\loigiai{
Xét hàng nghìn của số $2\,545\,000$ có chữ số là $5 \ge 5$, do đó khi làm tròn đến hàng chục nghìn ta cộng thêm $1$ vào chữ số hàng chục nghìn ($4 + 1 = 5$), ta được số $2\,550\,000$.\\
\textbf{Đáp án của câu hỏi là C}
}
\end{ex}

\textbf{Phần 2. Trắc nghiệm đúng sai}
\setcounterpageref{ex}{0}

\begin{ex} %[5D1H1]
Cho hai số tự nhiên $A = 26\,400\,500$ và $B = 30\,008\,021$. Xét tính đúng sai của các khẳng định sau:
\choiceTF
{\True Số $A$ được đọc là: Hai mươi sáu triệu bốn trăm nghìn năm trăm.}
{Trong số $B$, chữ số $8$ thuộc hàng trăm nghìn.}
{\True Số $A$ được viết thành tổng theo các hàng là $20\,000\,000 + 6\,000\,000 + 400\,000 + 500$.}
{\True Hiệu giữa giá trị của chữ số $3$ và giá trị của chữ số $8$ trong số $B$ là $29\,992\,000$.}
\loigiai{
\begin{itemize}
    \item Số $A = 26\,400\,500$ đọc là: Hai mươi sáu triệu bốn trăm nghìn năm trăm.
    \item Trong số $B = 30\,008\,021$, chữ số $8$ thuộc hàng nghìn (lớp nghìn), có giá trị là $8\,000$.
    \item Phân tích cấu tạo số $A$: $26\,400\,500 = 20\,000\,000 + 6\,000\,000 + 400\,000 + 500$.
    \item Trong số $B$: Chữ số $3$ thuộc hàng chục triệu có giá trị $30\,000\,000$; chữ số $8$ có giá trị $8\,000$. Hiệu giá trị là $30\,000\,000 - 8\,000 = 29\,992\,000$.
\end{itemize}
\begin{itemize}
    \item a) \textbf{Đúng}.
    \item b) \textbf{Sai}: Chữ số $8$ thuộc hàng nghìn.
    \item c) \textbf{Đúng}.
    \item d) \textbf{Đúng}.
\end{itemize}
}
\end{ex}

\begin{ex} %[5D1H1]
Số sản phẩm bán được của một cửa hàng trong bốn năm được thống kê như sau: Năm 2020 bán $2\,873$ sản phẩm; Năm 2021 bán $2\,837$ sản phẩm; Năm 2022 bán $3\,293$ sản phẩm; Năm 2023 bán $3\,018$ sản phẩm. Doanh thu một ngày là $2\,545\,000$ đồng. Xét tính đúng sai của các khẳng định sau:
\choiceTF
{\True Năm $2022$ cửa hàng bán được nhiều sản phẩm nhất.}
{Năm $2021$ cửa hàng bán được nhiều sản phẩm hơn năm $2020$.}
{\True Làm tròn doanh thu $2\,545\,000$ đồng đến hàng trăm nghìn ta được $2\,500\,000$ đồng.}
{Làm tròn doanh thu $2\,545\,000$ đồng đến hàng chục nghìn ta được $2\,540\,000$ đồng.}
\loigiai{
\begin{itemize}
    \item So sánh số sản phẩm: $2\,837 < 2\,873 < 3\,018 < 3\,293$. Vậy năm $2022$ bán nhiều sản phẩm nhất ($3\,293$ sản phẩm).
    \item Năm $2021$ bán $2\,837$ sản phẩm, năm $2020$ bán $2\,873$ sản phẩm. Vì $2\,837 < 2\,873$ nên năm $2021$ bán ít hơn năm $2020$.
    \item Chữ số hàng chục nghìn của $2\,545\,000$ là $4 < 5$, nên làm tròn đến hàng trăm nghìn được $2\,500\,000$ đồng.
    \item Chữ số hàng nghìn của $2\,545\,000$ là $5 \ge 5$, nên làm tròn đến hàng chục nghìn phải được $2\,550\,000$ đồng.
\end{itemize}
\begin{itemize}
    \item a) \textbf{Đúng}.
    \item b) \textbf{Sai}: Năm 2021 bán ít hơn năm 2020.
    \item c) \textbf{Đúng}.
    \item d) \textbf{Sai}: Kết quả làm tròn đến hàng chục nghìn là $2\,550\,000$ đồng.
\end{itemize}
}
\end{ex}

\textbf{Phần 3. Trả lời ngắn}
\setcounterpageref{ex}{0}

\begin{ex} %[5D1H1]
Tìm số thích hợp điền vào dấu hỏi chấm để hoàn thành đẳng thức sau:
\[
1\,730\,539 = 1\,000\,000 + ? + 30\,000 + 500 + 30 + 9.
\]
\shortans{700000}
\loigiai{
\textbf{Đáp số:} $700000$\\
Chữ số $7$ trong số $1\,730\,539$ ở hàng trăm nghìn, do đó có giá trị là $700\,000$.
}
\end{ex}

\begin{ex} %[5D1H1]
Ba bạn Việt, Nam và Rô-bốt đội $3$ chiếc mũ mang $3$ số chẵn liên tiếp. Rô-bốt nhìn thấy mũ của Nam ghi số $2\,028$ và mũ của Việt ghi số $2\,032$. Hỏi chiếc mũ của Rô-bốt ghi số nào?
\shortans{2030}
\loigiai{
\textbf{Đáp số:} $2030$\\
Ba số chẵn liên tiếp là $2\,028$; $2\,030$; $2\,032$. Vậy số trên mũ Rô-bốt là $2030$.
}
\end{ex}

\begin{ex} %[5D1H1]
Rô-bốt có thói quen viết các số biểu diễn ngày, tháng, năm liên tiếp nhau để được một số tự nhiên có nhiều chữ số. Ví dụ: ngày 30 tháng 4 năm 1975 viết thành $3\,041\,975$. Hỏi với ngày Nhà giáo Việt Nam năm 2024 (ngày 20 tháng 11 năm 2024), số tự nhiên Rô-bốt viết được là bao nhiêu?
\shortans{20112024}
\loigiai{
\textbf{Đáp số:} $20112024$\\
Ngày $20$, tháng $11$, năm $2024$ ghép liên tiếp thành số $20\,112\,024$.
}
\end{ex}

\begin{ex} %[5D1H1]
Trong số $20\,112\,024$, tổng giá trị của tất cả các chữ số $2$ có trong số đó bằng bao nhiêu?
\shortans{20002020}
\loigiai{
\textbf{Đáp số:} $20002020$\\
Các chữ số $2$ có giá trị lần lượt là:
\begin{itemize}
    \item Hàng chục triệu: $20\,000\,000$.
    \item Hàng nghìn: $2\,000$.
    \item Hàng chục: $20$.
\end{itemize}
Tổng giá trị của các chữ số $2$ là: $20\,000\,000 + 2\,000 + 20 = 20\,002\,020$.
}
\end{ex}

\begin{ex} %[5D1H1]
Doanh thu thực tế cả ngày của một cửa hàng là $2\,545\,000$ đồng. Hãy làm tròn doanh thu này đến hàng chục nghìn.
\shortans{2550000}
\loigiai{
\textbf{Đáp số:} $2550000$\\
Chữ số hàng nghìn là $5$, làm tròn lên hàng chục nghìn ta được $2\,550\,000$.
}
\end{ex}

\begin{ex} %[5D1V1]
Rô-bốt lập số $863\,749$ bằng các tấm thẻ: $8$, $6$, $3$, $7$, $4$, $9$. Rô-bốt đổi chỗ đúng $2$ tấm thẻ bất kỳ để nhận được một số lẻ lớn nhất có thể. Hỏi số lẻ lớn nhất đó là bao nhiêu?
\shortans{873649}
\loigiai{
\textbf{Đáp số:} $873649$\\
Số ban đầu là $863\,749$ (có chữ số tận cùng là $9$ nên là số lẻ).
\begin{itemize}
    \item Để số tạo thành là số lẻ lớn nhất, chữ số tận cùng phải là chữ số lẻ ($3$, $7$ hoặc $9$).
    \item Nếu đổi thẻ $9$ (hàng đơn vị) với thẻ $8$ (hàng trăm nghìn), ta được số $963\,748$, có tận cùng là $8$ (số chẵn, không thỏa mãn).
    \item Nếu đổi thẻ $9$ với $7$ ta được số $863\,947$ (lẻ). Nếu đổi thẻ $9$ với $3$ ta được $869\,743$ (lẻ).
    \item Nếu giữ thẻ $9$ ở hàng đơn vị và đổi chỗ thẻ $6$ (hàng chục nghìn) với thẻ $7$ (hàng trăm), ta được số $873\,649$ (lẻ).
    \item So sánh các số lẻ tạo được: $873\,649 > 869\,743 > 867\,349 > 863\,947 > 863\,749$.
\end{itemize}
Vậy số lẻ lớn nhất có thể nhận được là $873\,649$ (bằng cách đổi chỗ thẻ số $6$ và thẻ số $7$).
}
\end{ex}
